\documentclass[12pt,a4paper]{article}

% --- PAQUETES ---
\usepackage[utf8]{inputenc}
\usepackage[T1]{fontenc}
\usepackage[spanish,es-nodecimaldot]{babel}
\usepackage{amsmath, amssymb, amsthm, mathtools}
\usepackage{geometry}
\usepackage{xcolor}

% --- CONFIGURACIÓN ---
\geometry{left=25mm, right=25mm, top=25mm, bottom=25mm}
\newtheorem{teorema}{Teorema}
\newtheorem{corolario}{Corolario}
\newtheorem{proposicion}{Proposición}
\newtheorem{lema}{Lema}
\theoremstyle{definition}
\newtheorem{definicion}{Definición}
\newtheorem{ejemplo}{Ejemplo}
\theoremstyle{remark}
\newtheorem{observacion}{Observación}

% --- COMANDOS ---
\newcommand{\C}{\mathbb{C}}
\newcommand{\Poly}{\mathbb{P}[z]}
\newcommand{\norm}[1]{\left\| #1 \right\|_S}
\newcommand{\abs}[1]{\left| #1 \right|}
\newcommand{\ip}[2]{\langle #1, #2 \rangle_S}
\newcommand{\D}{\mathcal{D}}

\title{Análisis Espectral del Operador Multiplicación en Espacios de Sobolev}
\author{}
\date{}

\begin{document}

\maketitle
\section{Capítulo 1. Introducción y Preliminares 1 }
\begin{enumerate}
  


   
\item  Contexto y Motivación




1.2 Medidas de Borel, Regularidad y Medida de Lebesgue 



\item 1.3 Espacios de Polinomios y Ortogonalidad 

POLINOMIOS ORTOGONALES CLÁSICOS. 

Un apartado donde se explican el caso standard de polinomios ortogonales asociados a medidas. Aquí se introducen medidas de Borel, y la de Lebesgue, 
Se introduce el operador multiplicación por $z$ y se da el resultado general en el caso clásico con referencias de que son equivalentes: 
\begin{enumerate}
\item Acotación del operador. 
\item Acotación del soporte 
\item Acotación de ceros en la envoltura convexa.  
\end{enumerate}
\item  Espacios de Sobolev y Productos Discretos

\noindent Se definen los espacios de Sobolev discretos y resaltar las diferencias respecto del clásico. 

\begin{enumerate}
    \item El operador multiplicación no tiene porqué estar acotado aunque el soporte esté acotado REFERENCIAS Rodriguez comentar que daremos muchos ejemplos en el trabajo. 
    \item Ceros de polinomios los ceros no están en la envoltura convexa. Ejemplo de Althamer, etc....Problema de acotación de ceros. y Referencias y resultados de Pijeira, Rodríguez, Y 
\end{enumerate}

\item 1.5 Evaluaciones Puntuales Acotadas (BPE) y Puntos $\delta$-singulares  Cambiar título, sólo introducir  PUNTOS DE EVALUACIÓN CONTINUA ( LOS SINGULARES TODAVÍA NO) y los resultados nuestros 
sobre acotación de operador multiplicación y por tanto de ceros en este contexto. 
\end{enumerate}

1.6 Organización del Artículo NO 

\section{Capítulo 2 Análisis Espectral y Topología en Espacios de Sobolev 13 }

\begin{enumerate}
    \item  Codimensión y Subespacios en el Espacio de Hilbert
\noindent Introducir definiciones en espacios de polinomios $\mathbb{P}[z]$ con un producto escalar  y su complección $\mathcal{H}$ de: 

\begin{enumerate}
    \item Definir codimensión algebraica. Propiedades. Caracterizaciones como intesecciones de núcleos de funcionales acotados o no , evaluaciones en particular, (NO INTRODUCIR TODAVÍA PUNTOS SINGULARES NI LOS CONJUNTOS IN OUT), referencias. etc. 
\item Definir codimensión topológica, caracterizaciones como intersecciones de funcionales acotados, referencias, 
\end{enumerate}

\noindent IMPORTANTE: relaciones entre las dos codimensiones.  
\noindent Estudio de las codimensiones de núcleos de aplicaciones lineales, en particular de evaluaciones. 

    \item La Dicotomía Topológica
\item El Operador de Multiplicación y los Números de Gelfand

\noindent Introducir primero los números de Gelfand que son los conocidos, para operadores acotados en espacios de Hilbert, propiedades, etc...
\item Un Nuevo Índice de Acotación Algebraica: El Índice Qk
\noindent Introducir el $Q_k$ sus propiedades. 

\noindent Relación entre los índices de Gelfand y los $Q_k$'s para cualquier operador acotado. 
\end{enumerate}

\section{Valores singulares y operador multiplicación }

\noindent Este capítulo se centrará en el teorema principal de la estabilización de la cadena. Se definen los puntos singulares, los in y los out, y el teorema principal. 

\section{ Análisis Matricial y Comportamiento Asintótico}
3 Análisis Matricial y Comportamiento Asintótico

3.1 Representación Matricial y Secciones Truncadas


3.2 Valores Singulares y Operadores Truncados . . . . 


3.3 Estabilización Algebraica de los Valores Singulares 


3.4 Conteo de Ceros Excepcionales: la Desigualdad deWeyl-Horn 


3.5 Relación entre los Valores Singulares y el Índice Qk 


3.6 Matrices de Momentos de Sobolev y de Dirac


4 Localización y Acotación Asintótica de los Ceros 63
4.1 Los Ceros como Autovalores 


4.2 Acotación uniforme vía cuasi-ortogonalidad 


4.3 La cadena estructural: ?  Equivalencia y comparación de las dos vías 



\section{Correcciones Introducción}

\end{document}